\documentclass[10pt,a4paper]{amsart}
\usepackage[margin=19mm]{geometry}
\usepackage{amsmath,amssymb,mathtools}
\usepackage[scr=rsfs,cal=euler]{mathalfa}
\usepackage{xfrac}
\usepackage[unicode,hidelinks,bookmarks=false]{hyperref}
\newcommand{\ie}{i.e.\ }
\newcommand{\cf}{cf.\ }
\newcommand{\resp}{respectively\ }
\newcommand{\Subsection}{\subsection}
\providecommand{\nobreakdash}{\nobreak\hbox{-}}
\setlength{\emergencystretch}{2em}
\multlinegap0pt
\usepackage{amssymb}
\newtheorem{lemma}{Lemma}
\title{Locally Conformally Flat Manifolds with Positive Scalar Curvature: Kleinian Groups, Moduli Spaces, and Euclidean Rigidity}
\author{Jialong Deng}
\date{Source: arXiv:2609.02267v1 (CC BY 4.0)}
\begin{document}
\maketitle

\noindent\emph{Source and licence.} Locally Conformally Flat Manifolds with Positive Scalar Curvature: Kleinian Groups, Moduli Spaces, and Euclidean Rigidity. Version arXiv:2609.02267v1 (CC BY 4.0), distributed by arXiv under the Creative Commons Attribution 4.0 International licence, http://creativecommons.org/licenses/by/4.0/, as recorded in the arXiv licence metadata for this exact version and shown on the abstract page https://arxiv.org/abs/2609.02267v1. Full licence notice: \texttt{licenses/arxiv-2609.02267-CC-BY-4.0.txt}.\par\smallskip

\noindent\textit{Editorial scope.} Counted unit: the article's lemma lem:cylinder-moduli (source0.tex lines 3114--3116). The article's complete original proof of it is delivered with it (source0.tex lines 3118--3437, one \texttt{proof} environment of 320 lines). No mathematics is written, repaired or re-derived: the statement and the proof are the article's own lines, in the article's own notation, and no retained line is substituted or re-flowed.  No further source-local context is needed: every cross-reference of every retained line resolves inside the two delivered spans.  Nothing was omitted from any retained span, so the package carries no omission record.  No retained line contains a citation, so the wrapper carries no bibliography and no bibliography entry.  No retained line contains a figure environment, an image-inclusion command or drawing code, so no image is delivered and the package declares no asset dependency.  Editorial formatting only, in addition: an AMS \texttt{amsart} preamble that restates the article's own declarations and macros used by the retained lines, namely the environments \texttt{lemma} on separate counters, together with the standard packages those lines need.  The retained environments print in this document as Lemma 1 in order of appearance, whereas the article prints the same items with the numbering of its own full document.  The complete native source of the article as distributed in the arXiv e-print for this exact version is bundled unchanged as \texttt{sources/209186/source0.tex}.
\begin{lemma}\label{lem:cylinder-moduli}
\(\mathcal{M}^{CC}(S^2 \times S^1)\) is homeomorphic to \((0,\infty) \times [0,\pi]\).
\end{lemma}

\begin{proof}
We construct continuous maps
\(Q\colon(0,\infty)\times[0,\pi]\to  \mathcal{M}^{CC}(S^2\times S^1)\)
and
\(\overline P\colon\mathcal{M}^{CC}(S^2\times S^1)\to (0,\infty)\times[0,\pi]\)
and prove that they are inverse homeomorphisms.

Fix once and for all the universal covering
\[
\varpi\colon S^2\times \mathbb R\longrightarrow S^2\times S^1,
\qquad
\varpi(x,s)=(x,[s]),
\]
with deck generator
\[
\tau\colon S^2\times \mathbb R\longrightarrow S^2\times \mathbb R,
\qquad
\tau(x,s)=(x,s+1).
\]
Thus \(\mathrm{Deck}(\varpi)=\langle \tau\rangle\cong \mathbb Z\).

Let \(h\in \mathrm{Met}^{CC}(S^2\times S^1)\), and write
$\widetilde h:=\varpi^*h$. By part~(i) of the local-to-global cylinder
observation, there is an isometry
\(J_h\colon(S^2\times\mathbb R,\widetilde h)
\to (S^2\times\mathbb R,g_{\mathrm{cyl}})\).
Set
\(\gamma:=J_h\circ\tau\circ J_h^{-1}
\in\mathrm{Isom}(S^2\times\mathbb R,g_{\mathrm{cyl}})\).
Because \(\gamma\) is conjugate to \(\tau\), the action of \(\langle\gamma\rangle\) is free and \(\gamma\) has infinite order.

The Ricci endomorphism of the round cylinder has eigendistributions $TS^2$
and $T\mathbb R$, with the distinct eigenvalues $1$ and $0$. Hence every
isometry preserves both distributions. Writing an isometry as
$\Psi=(\Psi_1,\Psi_2)$, this gives $\partial_t\Psi_1=0$ and
$d_x\Psi_2=0$. Since both factors are connected, it follows that
$\Psi(x,t)=(A(x),b(t))$. Bijectivity of $\Psi$ makes $A$ and $b$
diffeomorphisms, while $\Psi^*g_{\mathrm{cyl}}=g_{\mathrm{cyl}}$ gives
$A^*g_{st}=g_{st}$ and $b^*dt^2=dt^2$. Consequently,
\[
\mathrm{Isom}(S^2\times \mathbb R, g_{\mathrm{cyl}})
\cong \mathrm{Isom}(S^2, g_{st}) \times \mathrm{Isom}(\mathbb{R}, dt^2)
\cong \mathrm{O}(3) \times (\mathbb{R} \rtimes \{\pm 1\}).
\]
Thus
\(
\gamma(x,t)=(Ax,\delta t+L)
\)
for some $A\in\mathrm{O}(3)$, $\delta\in\{\pm1\}$, and $L\in\mathbb R$.


We claim that \(\delta=1\). If \(\delta=-1\), then
$\gamma^2(x,t)=(A^2x,t)$.
Since \(A^2\in \mathrm{SO}(3)\), there exists \(x_0\in S^2\) such that
\(A^2x_0=x_0\). Hence $\gamma^2(x_0,t)=(x_0,t)$ for all $t\in\mathbb R$.
Since \(\gamma\) has infinite order, \(\gamma^2\neq\mathrm{id}\), contradicting
the freeness of the \(\langle\gamma\rangle\)-action. Thus \(\delta=1\).

We next claim that \(A\in \mathrm{SO}(3)\). The manifold
\(S^2\times S^1\) is oriented, and \(\tau\) preserves the lifted orientation
on \(S^2\times\mathbb R\). Since orientation preservation is invariant under
conjugation, the conjugate $\gamma=J_h\tau J_h^{-1}$ also preserves the product orientation.
The translation \(t\mapsto t+L\) preserves the orientation of $\mathbb R$;
hence $A$ preserves the orientation of $S^2$. Therefore
$A\in\mathrm{SO}(3)$.

We have \(L\neq 0\). Indeed, if \(L=0\), then $\gamma(x,t)=(Ax,t).$
Since \(A\in \mathrm{SO}(3)\), there exists \(x_0\in S^2\) such that \(Ax_0=x_0\). Hence
$\gamma(x_0,t)=(x_0,t)$ for all $t\in \mathbb R$, again contradicting freeness. Thus \(L\neq 0\).

Finally, after postcomposing $J_h$ with the reflection
$(x,t)\mapsto(x,-t)$ if necessary, we may assume $L>0$. After this
normalization, the parameters do not depend on the choices. Indeed, if
$J_h'=\Psi\circ J_h$, where
$\Psi(x,t)=(Bx,\varepsilon t+c)$, then
\((\Psi\gamma\Psi^{-1})(x,t)
=(BAB^{-1}x,t+\varepsilon L)\).
After restoring the convention that the translation length is positive,
$L$ is unchanged and $A$ is replaced by an $\mathrm O(3)$-conjugate. If the
deck generator $\tau$ is replaced by $\tau^{-1}$, then the holonomy element becomes
$\gamma^{-1}(x,t)=(A^{-1}x,t-L)$; conjugating by $(x,t)\mapsto(x,-t)$ to
restore the normalization gives $(A^{-1}x,t+L)$. Thus neither operation
changes the positive translation length $L$ or the rotation angle
$\theta\in[0,\pi]$ of $A$. Hence
$h$ determines the well-defined parameters $(L,\theta)$, and its holonomy
admits the normalized representative
\[
\gamma_{L,A}(x,t):=(Ax,t+L),
\qquad
A\in \mathrm{SO}(3),\quad L>0.
\]

We now construct the continuous map
$Q\colon (0,\infty) \times [0,\pi] \to \mathcal{M}^{CC}(S^2 \times S^1)$.
For \((L,\theta)\in (0,\infty)\times[0,\pi]\), let
\(
\gamma_{L,\theta}(x,t):=(R_\theta x,t+L),
\)
where \(R_\theta\in\mathrm{SO}(3)\) denotes the standard rotation by angle \(\theta\) about the \(x_3\)-axis. Let
$(S^2\times \mathbb R)/\langle \gamma_{L,\theta}\rangle$
carry the quotient metric \(\bar g_{L,\theta}\).

To identify this quotient with \(S^2\times S^1\), define
\[
\widetilde F_{L,\theta}\colon S^2\times \mathbb R\longrightarrow S^2\times \mathbb R,
\qquad
\widetilde F_{L,\theta}(x,s)=(R_{s\theta}x,Ls).
\]
This is a diffeomorphism, with inverse
\(\widetilde F_{L,\theta}^{-1}(y,t)=\bigl(R_{-(t/L)\theta}y,t/L\bigr)\).
Moreover,
\(\widetilde F_{L,\theta}\circ \tau=\gamma_{L,\theta}\circ \widetilde F_{L,\theta}\).
Hence \(\widetilde F_{L,\theta}\) descends to a diffeomorphism
\(F_{L,\theta}\colon S^2\times S^1\to
(S^2\times \mathbb R)/\langle \gamma_{L,\theta}\rangle\).
Define
\(h_{L,\theta}:=F_{L,\theta}^*\bar g_{L,\theta}\in \mathrm{Met}^{CC}(S^2\times S^1)\).

On the fixed universal cover one has
\(\varpi^*h_{L,\theta}=\widetilde F_{L,\theta}^*g_{\mathrm{cyl}}\).
Let $Z$ be the Killing field on $S^2$ generating $a\mapsto R_a$. For
$(V,a),(W,b)\in T_xS^2\oplus\mathbb R\partial_s$, direct differentiation
gives
\[
(\varpi^*h_{L,\theta})\bigl((V,a),(W,b)\bigr)
=g_{st}(V+a\theta Z_x,W+b\theta Z_x)+L^2ab.
\]
This tensor is $\tau$-invariant and depends smoothly on $(L,\theta)$,
including at $\theta=0$ and $\theta=\pi$. Hence the map
$(L,\theta)\longmapsto h_{L,\theta}$
is continuous in the \(C^\infty\)-topology. Thus there is a continuous map
\[
Q\colon(0,\infty)\times[0,\pi]\longrightarrow \mathcal{M}^{CC}(S^2\times S^1),
\qquad
Q(L,\theta):=[h_{L,\theta}].
\]


We now construct the map \(\overline P\).
Let \(h\in \mathrm{Met}^{CC}(S^2\times S^1)\), and let
$\widetilde h:=\varpi^*h$.
Define
\[
f_h\colon S^2\times \mathbb R\longrightarrow \mathbb R,
\qquad
f_h(p):=d_{\widetilde h}(p,\tau p).
\]
Since \(\tau\) is an isometry of \((S^2\times \mathbb R,\widetilde h)\), one has
$f_h(\tau p)=f_h(p),$
so \(f_h\) descends to a continuous function on the compact quotient \(S^2\times S^1\), still denoted \(f_h\).

Set
\[
L(h):=\min_{S^2\times S^1} f_h,
\qquad
D(h):=\max_{S^2\times S^1} f_h,
\qquad
\Theta(h):=\sqrt{D(h)^2-L(h)^2}.
\]
Since $\tau$ acts freely, $f_h(p)>0$ for every $p$. Since $f_h$ is continuous
on the compact quotient, it attains a positive minimum. Hence $L(h)>0$.

We now compute these invariants from the normalized holonomy. Let $J_h$ be the
isometry above satisfying $J_h\tau J_h^{-1}=\gamma_{L,A}$, where
$A\in\mathrm{SO}(3)$ has rotation angle $\theta\in[0,\pi]$ and $L>0$.
Because isometries preserve the distance function, the displacement $f_h(p)$
transforms as
\[
f_h(p)
=d_{g_{\mathrm{cyl}}}(J_h(p),J_h(\tau p))
=d_{g_{\mathrm{cyl}}}((x,t),\gamma_{L,A}(x,t)),
\]
where $(x,t):=J_h(p)$. Thus we obtain
\[
f_h(p)^2 = d_{g_{st}}(x, Ax)^2 + d_{\mathbb{R}}(t, t+L)^2 = d_{g_{st}}(x, Ax)^2 + L^2.
\]
Let $v$ be a unit vector on the rotation axis of $A$, chosen arbitrarily if
$\theta=0$, and put $a:=\langle x,v\rangle$. Then
\(\langle x,Ax\rangle=a^2+(1-a^2)\cos\theta\),
so $d_{g_{st}}(x,Ax)$ ranges over $[0,\theta]$. If $\theta>0$, its minimum
is attained on the rotation axis and its maximum on the orthogonal equator;
if $\theta=0$, it vanishes everywhere. It follows that
\(L(h)^2 = L^2\) and \(D(h)^2 = \theta^2 + L^2\).
Therefore
\(L(h)=L\), \(D(h)=\sqrt{L^2+\theta^2}\), and \(\Theta(h)=\theta\).

Thus
\[
P\colon\mathrm{Met}^{CC}(S^2\times S^1)\longrightarrow (0,\infty)\times[0,\pi],
\qquad
P(h):=\bigl(L(h),\Theta(h)\bigr)
\]
is well defined.

We next prove that \(P\) is continuous. Because the $C^\infty$-topology on
$\mathrm{Met}^{CC}(S^2\times S^1)$ is metrizable, it suffices to prove
sequential continuity. Let \(h_i\to h\) in the \(C^\infty\)-topology. In
particular, \(h_i\to h\) in \(C^0\). Since \(h\) is positive definite on the compact manifold \(S^2\times S^1\), after discarding finitely many terms and renumbering, we may choose
$\varepsilon_i\in(0,1),$ $\varepsilon_i\to 0,$ such that
\[
(1-\varepsilon_i)h\le h_i\le (1+\varepsilon_i)h
\]
as bilinear forms on \(T(S^2\times S^1)\). Pulling back by \(\varpi\), we obtain
\[
(1-\varepsilon_i)\widetilde h\le \widetilde h_i\le (1+\varepsilon_i)\widetilde h
\]
on \(T(S^2\times \mathbb R)\), where \(\widetilde h_i:=\varpi^*h_i\).

For every piecewise \(C^1\) curve \(c\) in \(S^2\times \mathbb R\),
\[
\sqrt{1-\varepsilon_i}\,\ell_{\widetilde h}(c)
\le
\ell_{\widetilde h_i}(c)
\le
\sqrt{1+\varepsilon_i}\,\ell_{\widetilde h}(c).
\]
Taking infima over curves joining fixed endpoints yields
\[
\sqrt{1-\varepsilon_i}\,d_{\widetilde h}
\le
d_{\widetilde h_i}
\le
\sqrt{1+\varepsilon_i}\,d_{\widetilde h}.
\]
Hence
\[
\sqrt{1-\varepsilon_i}\,f_h\le f_{h_i}\le \sqrt{1+\varepsilon_i}\,f_h
\]
pointwise on \(S^2\times S^1\). If
\(\eta_i:=\max\!\left\{\sqrt{1+\varepsilon_i}-1,\ 1-\sqrt{1-\varepsilon_i}\right\}\),
then \(\eta_i\to 0\) and
\(\|f_{h_i}-f_h\|_{C^0(S^2\times S^1)}\le \eta_i\,D(h)\to 0\).
Therefore
\(L(h_i)\to L(h)\), \(D(h_i)\to D(h)\), and \(\Theta(h_i)\to \Theta(h)\).
Thus \(P\) is continuous.

We next establish that $P$ is invariant under pullback by
$\mathrm{Diff}(S^2\times S^1)$. Let $\psi\in\mathrm{Diff}(S^2\times S^1)$,
and let $\widetilde\psi\colon S^2\times\mathbb R\to S^2\times\mathbb R$ be a
lift. Conjugation by $\widetilde\psi$ induces an automorphism of
$\mathrm{Deck}(\varpi)\cong\mathbb Z$. Hence there exists
$\varsigma\in\{\pm1\}$ such that
\(\widetilde\psi\circ\tau\circ\widetilde\psi^{-1}=\tau^\varsigma\).
By naturality of pullback,
\(\varpi^*(\psi^*h)=\widetilde\psi^*(\varpi^*h)
=\widetilde\psi^*\widetilde h\).
Therefore, for every $p\in S^2\times\mathbb R$,
\[
\begin{aligned}
f_{\psi^*h}(p)
&=d_{\widetilde\psi^*\widetilde h}(p,\tau p)\\
&=d_{\widetilde h}(\widetilde\psi p,\widetilde\psi\tau p)\\
&=d_{\widetilde h}(\widetilde\psi p,\tau^\varsigma\widetilde\psi p).
\end{aligned}
\]
If $\varsigma=1$, the last expression equals $f_h(\widetilde\psi p)$. If
$\varsigma=-1$, the symmetry of the distance function and the fact that
$\tau$ is an isometry of $\widetilde h$ give
\(d_{\widetilde h}(\widetilde\psi p,\tau^{-1}\widetilde\psi p)
=d_{\widetilde h}(\tau\widetilde\psi p,\widetilde\psi p)
=f_h(\widetilde\psi p)\).
Thus $f_{\psi^*h}=f_h\circ\widetilde\psi$ in both cases. The two functions
have the same range, and hence
\(L(\psi^*h)=L(h)\), \(D(\psi^*h)=D(h)\), and \(\Theta(\psi^*h)=\Theta(h)\).
It follows that $P$ is $\mathrm{Diff}(S^2\times S^1)$-invariant. By the
universal property of the quotient, $P$ descends to a unique continuous map
\(\overline P\colon\mathcal{M}^{CC}(S^2\times S^1)
\to (0,\infty)\times[0,\pi]\).

We conclude by showing that $Q$ and $\overline P$ are inverse homeomorphisms.
For any $(L,\theta)\in(0,\infty)\times[0,\pi]$, the construction of the
model metric $h_{L,\theta}$ provides an isometry
\(\widetilde F_{L,\theta}\colon
(S^2\times\mathbb R,\varpi^*h_{L,\theta})
\to (S^2\times\mathbb R,g_{\mathrm{cyl}})\)
that conjugates $\tau$ to
$\gamma_{L,\theta}(x,t)=(R_\theta x,t+L)$. For
$p\in S^2\times\mathbb R$, write
\(
(x,t):=\widetilde F_{L,\theta}(p).
\)
Then
\(f_{h_{L,\theta}}(p)
=\sqrt{d_{g_{st}}(x,R_\theta x)^2+L^2}\).
By the displacement computation above,
$d_{g_{st}}(x,R_\theta x)$ ranges over $[0,\theta]$; in particular, this
statement also covers $\theta=0$. It follows that
$L(h_{L,\theta})=L$ and $\Theta(h_{L,\theta})=\theta$. Thus,
\(\overline P(Q(L, \theta)) = \overline P([h_{L, \theta}]) = (L, \theta)\).
Conversely, let $h\in\mathrm{Met}^{CC}(S^2\times S^1)$ have invariants
$P(h)=(L,\theta)$. By the holonomy normalization established above, there is
an isometry
\(J_h\colon(S^2\times\mathbb R,\varpi^*h)
\to (S^2\times\mathbb R,g_{\mathrm{cyl}})\)
such that $J_h\tau J_h^{-1}=\gamma_{L,A}$, where $A\in\mathrm{SO}(3)$
has rotation angle $\theta$. Choose $C\in\mathrm{SO}(3)$ such that
$CAC^{-1}=R_\theta$, and set $K_C(x,t):=(Cx,t)$. Then
\((K_C\circ J_h)\tau(K_C\circ J_h)^{-1}
=K_C\gamma_{L,A}K_C^{-1}=\gamma_{L,\theta}\).
Hence $K_C\circ J_h$ descends to an isometry
\[
\overline K\colon
(S^2\times S^1,h)
\longrightarrow
\bigl((S^2\times\mathbb R)/\langle\gamma_{L,\theta}\rangle,
\bar g_{L,\theta}\bigr).
\]
Composing with $F_{L,\theta}^{-1}$ gives an isometry
\(F_{L,\theta}^{-1}\circ\overline K\colon
(S^2\times S^1,h)
\to (S^2\times S^1,h_{L,\theta})\).
Consequently,
\(Q(\overline P([h]))=Q(L,\theta)=[h_{L,\theta}]=[h]\).

The identities $\overline P\circ Q=\mathrm{id}$ and
$Q\circ\overline P=\mathrm{id}$, together with the continuity of $\overline P$
and $Q$,
show that $\overline P$ and $Q$ are inverse homeomorphisms. Thus
\(\mathcal{M}^{CC}(S^2\times S^1)\cong(0,\infty)\times[0,\pi]\).
\end{proof}

\end{document}
